\section{Matched tests of the final fusion readout}
\label{sec:design-analysis}
The official result compares complete systems. Here we isolate readout choices while holding the final RGB and event producers, temporal refiners, supplied regions, and causal histories fixed. We use the same 460 corrected transfer queries throughout. These recordings informed earlier development; the new comparisons quantify conditional effects on that cohort.

\subsection{Simple combinations of identical branch outputs}
We first replace the released fusion head with equal mixtures of the two final branch predictions, combining them in phase, inverse TTC, or seconds. Each mixture receives the same branch outputs and temporal support. All emit finite TTC under their specified numerical policies. Table~\ref{tab:revision-fixed} shows that learned fusion improves mean error over these fixed combinations. Its comparison with the strong RGB branch remains uncertain (Section~\ref{sec:late-transfer}).
\begin{table}[htbp]
\centering\small\begin{tabular}{lrrr}
\toprule
Frozen readout & MAE (s) & Median (s) & p95 (s) \\
\midrule
Released fusion & 0.813 & 0.446 & 2.742 \\
RGB context & 0.851 & 0.390 & 3.038 \\
Event context & 1.376 & 0.844 & 3.707 \\
Equal phase mixture & 0.997 & 0.519 & 2.759 \\
Equal inverse-TTC mixture & 1.007 & 0.518 & 2.762 \\
Equal TTC mixture & 1.050 & 0.614 & 3.192 \\
\bottomrule
\end{tabular}

\caption{Frozen final-branch readouts on the identical 460 queries. Equal mixtures have no fitted parameters; the released head has 1,955. This comparison holds upstream information fixed but does not match readout capacity.}
\label{tab:revision-fixed}
\end{table}

\subsection{Equal-capacity heads with matched fitting}
To distinguish event access from added correction capacity, we fit three heads: the original phase readout; an identical head whose ten event fields are zero during fitting and evaluation; and a residual readout expressed in inverse TTC. All contain 1,955 parameters, start with a zero final affine map, and retain the RGB anchor. The zero-event head still receives RGB timing and support. Upstream weights are frozen throughout.

Each arm uses the same 88,744 source-domain out-of-fold feature rows from forty eAP recordings, the same seeded minibatch sequence, and 12,044 AdamW updates at batch size 128. We retain the original warmup/cosine schedule, weight decay, clipping, and phase-error plus quantile objective. Seeds 7, 13, and 21 vary only the correction-head fit. EvTTC queries and targets are excluded from fitting and checkpoint selection; every arm uses its final update. The inverse-TTC arm replaces phase inputs with normalized rates and predicts an RGB rate correction before conversion to the common phase loss. Appendix~\ref{app:consolidation-reproduction} specifies the transforms and archived source bindings.

We also fit a single scalar phase-mixture weight on the source training rows alone. Searching the specified 101-point grid gives event weight 0.31. This is a train-calibrated simple baseline, evaluated once on EvTTC.
\begin{table}[htbp]
\centering\small\begin{tabular}{lrr}
\toprule
Readout & MAE (s) & Relative error (\%) \\
\midrule
Phase correction, real events & $0.847\pm0.038$ & $30.92\pm2.05$ \\
Phase correction, zero events & $0.838\pm0.021$ & $29.12\pm1.16$ \\
Inverse-TTC correction, real events & $0.847\pm0.041$ & $30.80\pm2.14$ \\
Source-calibrated scalar phase mixture & 0.996 & 36.23 \\
\bottomrule
\end{tabular}

\caption{Matched final-head fitting. Values are the mean and sample standard deviation of the three seed-level scores; predictions are not ensembled. The scalar mixture has one fitted coefficient and no seed-dependent optimization.}
\label{tab:revision-heads}
\end{table}

\subsection{What the controlled comparison establishes}
Real events improve MAE over the zero-event control for one seed and worsen it for two. The mean zero-minus-real difference is $-0.0094$\,s, with a descriptive recording-bootstrap interval spanning zero. Inverse-TTC and phase corrections are also indistinguishable at this resolution. We average paired losses over the three fixed seeds before resampling complete recordings, preserving query correspondence; seeds do not multiply the number of independent recordings.
\begin{table}[htbp]
\centering\small\begin{tabular}{lrr}
\toprule
Paired contrast & $\Delta$MAE (s) & 95\% interval (s) \\
\midrule
Zero events minus real events & -0.0094 & $[-0.0780,+0.0630]$ \\
Inverse TTC minus phase & -0.0003 & $[-0.0061,+0.0058]$ \\
Calibrated mixture minus phase & +0.1492 & $[+0.0064,+0.3762]$ \\
\bottomrule
\end{tabular}

\caption{Paired differences after averaging query losses over fixed head seeds; 10,000 complete-recording resamples across 32 recordings. These intervals describe the inspected cohort and condition on the shared upstream models.}
\label{tab:revision-paired}
\end{table}

The learned readouts improve on the calibrated scalar mixture, but the tests give no consistent incremental benefit from event access and no observed advantage of phase over the tested inverse-TTC parameterization. A useful RGB correction, rather than event complementarity alone, can therefore account for the readout's gain on this cohort. The prior positive event control used a different adapted model; Appendix~\ref{app:design-history} preserves that result without transferring its attribution to the submitted system. A matched hidden-feature fusion model and genuinely unused recordings remain necessary to compare broader fusion designs.
\FloatBarrier
